\documentclass[11pt]{article}

\usepackage[margin=1in]{geometry}
\usepackage{amsmath,amssymb,amsthm}
\usepackage{booktabs}
\usepackage{graphicx}
\usepackage{hyperref}
\hypersetup{colorlinks=true, linkcolor=blue, citecolor=blue, urlcolor=blue}

\newtheorem{theorem}{Theorem}
\newtheorem{proposition}{Proposition}
\newtheorem{definition}{Definition}
\newtheorem{remark}{Remark}

\newcommand{\E}{\mathbb{E}}
\newcommand{\Prob}{\mathbb{P}}
\newcommand{\R}{\mathbb{R}}

\title{Sharp Composition Bounds for AI Guardrails Under Unknown Dependence}
\author{Pranav Bhave}
\date{\today}

\begin{document}
\maketitle

\begin{abstract}
AI systems are increasingly protected by stacks of guardrails, yet evaluations
often report only singleton failure rates for each guardrail. Those marginals do
not identify the failure probability of a composed stack unless the dependence
structure among guardrail failures is measured or assumed. We formalize
composed guardrail evaluation as a finite partial-identification problem over
binary failure events. Given singleton marginals, optional pairwise or linear
side constraints, and a declared Boolean composition event, the method computes
sharp lower and upper bounds by linear programming over the atom simplex. Here
``sharp'' means the endpoints are attainable under the declared constraints; it
does not mean the resulting interval is narrow or sufficient for deployment.
The framework reports identified-set width, product-coupling regret, endpoint
witness distributions, sample-complexity guidance, and explicit scaling limits.
\end{abstract}

\section{Introduction}
\input{sections/01_introduction.tex}

\section{Problem Setup}
Let $G_1,\ldots,G_m$ denote a finite collection of guardrails. For each
guardrail define a binary failure event
\[
  Z_i \in \{0,1\}, \qquad Z_i = 1 \text{ means guardrail failure or unsafe pass.}
\]
The unknown object is the atom distribution
\[
  \pi(z) = \Prob(Z=z), \qquad z \in \{0,1\}^m.
\]
The atom ordering used by the implementation is deterministic and little
endian, but the mathematical claims are invariant to relabeling.

Evaluation evidence may include singleton failure marginals
\[
  p_i = \Prob(Z_i=1)
\]
and optional pairwise overlap constraints
\[
  q_{ij} = \Prob(Z_i=1, Z_j=1),
\]
possibly as intervals. More generally, the implementation accepts linear
constraints of the form $a^\top \pi \le b$, $a^\top \pi \ge b$, or
$a^\top \pi=b$.

A composed safety question is a declared event $\phi:\{0,1\}^m\to\{0,1\}$.
Examples include at least one selected guardrail failing, all selected
guardrails failing, or a custom event over the finite atom space. Because
$\E_\pi[\phi(Z)]$ is linear in $\pi$, composition bounds are linear programs
over a finite simplex.

The distinction between mathematical and empirical assumptions matters. The LP
can verify what follows from supplied rates and constraints. It cannot verify
that those rates were measured on a representative dataset, that labels are
correct, or that future deployment prompts follow the evaluated population.


\section{Sharp Finite-Atom Bounds}
Let $\mathcal{F}$ be the feasible set of atom distributions satisfying the
simplex constraints and all declared evidence:
\[
  \mathcal{F} = \{\pi \in \R^{2^m}: \pi_z \ge 0,\ \sum_z \pi_z = 1,\ A\pi \le b,\ C\pi=d\}.
\]
For a query $\phi$, define
\[
  L_\phi = \inf_{\pi\in\mathcal{F}}\E_\pi[\phi(Z)],
  \qquad
  U_\phi = \sup_{\pi\in\mathcal{F}}\E_\pi[\phi(Z)].
\]

\begin{theorem}[Finite sharpness]
If $\mathcal{F}$ is nonempty, then $L_\phi$ and $U_\phi$ are attained by
feasible atom distributions. Every probability in $[L_\phi,U_\phi]$ is
compatible with the declared linear constraints after mixing endpoint
witnesses.
\end{theorem}

\begin{proof}
The feasible set is a closed and bounded polytope in finite-dimensional
Euclidean space, hence compact. The query functional is linear and therefore
continuous, so it attains its minimum and maximum on $\mathcal{F}$. Convex
mixtures of feasible distributions remain feasible and interpolate the query
value between endpoint optima.
\end{proof}

Classical Frechet-Hoeffding bounds are recovered when the only evidence is
singleton marginals. For an intersection event,
\[
  \Prob(\cap_i \{Z_i=1\}) \in
  \left[\max\left\{0,\sum_i p_i-(m-1)\right\},\ \min_i p_i\right],
\]
and for a union event,
\[
  \Prob(\cup_i \{Z_i=1\}) \in
  \left[\max_i p_i,\ \min\left\{1,\sum_i p_i\right\}\right].
\]
These closed forms are not exponential to compute. Exponential scaling enters
when the general atom-LP path is used for arbitrary Boolean events or side
constraints: the explicit LP has $2^m$ atom variables.


\section{Diagnostics and Witnesses}
The primary output is the interval $[L_\phi,U_\phi]$. We report its width
\[
  W_\phi = U_\phi - L_\phi
\]
as an identified-set diagnostic. Large width means the supplied evidence leaves
substantial composition risk unidentified.

When a selected or observed event risk $r$ is available, we report its position
inside the interval:
\[
  \operatorname{pos}(r) = \frac{r-L_\phi}{U_\phi-L_\phi},
\]
with the degenerate case handled separately. This is descriptive, not a safety
claim.

We also compute the product-coupling baseline $r_{\mathrm{prod}}$ obtained by
treating singleton failures as independent Bernoulli variables. The
independence regret
\[
  r-r_{\mathrm{prod}}
\]
is a signed diagnostic for how far an observed or selected risk lies from that
explicit assumption. The product-coupling baseline is useful as a comparison
point, but it is not the default model of reality.

Every reported sharp interval should include endpoint witness distributions
$\pi_L$ and $\pi_U$. These witnesses satisfy the declared constraints and attain
the lower and upper query values. They make the mathematical result
recomputable and challengeable. They do not prove that the empirical inputs are
representative or that either endpoint is realized in deployment.


\section{Sample Complexity and Scaling}
The kernel consumes Bernoulli rates. Estimating those rates from data is a
separate statistical problem. For one Bernoulli rate estimated from $n$ samples,
Hoeffding's inequality gives
\[
  \Prob(|\hat p-p| > \epsilon) \le 2\exp(-2n\epsilon^2).
\]
Thus a two-sided radius at error probability $\delta$ is
\[
  \epsilon(n,\delta) = \sqrt{\frac{\log(2/\delta)}{2n}},
\]
and the corresponding sample size for radius $\epsilon$ is
\[
  n \ge \frac{\log(2/\delta)}{2\epsilon^2}.
\]

For $k$ simultaneously reported Bernoulli rates, a union bound replaces
$\delta$ by $\delta/k$, yielding
\[
  \epsilon(n,k,\delta) =
  \sqrt{\frac{\log(2k/\delta)}{2n}}.
\]
For $m$ guardrails, the number of singleton plus pairwise rates is
$m + m(m-1)/2$.

Computational scaling has a different source. Marginal-only AND and OR bounds
have closed forms. The general finite-atom LP enumerates all $2^m$ atoms, so it
is appropriate for small guardrail stacks, theorem validation, and auditable
examples. Scaling beyond that requires structure, decomposition, or a different
optimization layer.


\section{Worked Benchmark Example}
The worked example uses harmful-prompt benchmark rows and two guardrails:
LlamaGuard and a deterministic keyword blocker. The failure convention is
fixed: $Z_i=1$ means guardrail $i$ did not block an unsafe prompt. For a stack
that blocks if any guardrail blocks, the unsafe-pass event is
\[
  \phi_{\mathrm{pass}}(Z) = \mathbf{1}\{Z_1=1,\ldots,Z_m=1\}.
\]
We also report
\[
  \phi_{\mathrm{any}}(Z) = \mathbf{1}\{\exists i: Z_i=1\},
\]
which captures whether at least one guardrail failed on the prompt.

The deterministic keyword blocker is a fixed function of prompt text. On a
given dataset its overlap with LlamaGuard failures is exactly computable from
the realized rows, although generalization beyond that dataset remains a
statistical question. Continuous integration uses a mocked LlamaGuard adapter
only to validate the pipeline. Paper figures derived from a real model run must
record the model identifier, immutable revision, dataset hash, adapter
configuration, and run manifest.

The example reports both marginal-only intervals and pairwise-tightened
plug-in intervals. If the marginal-only interval is wide, that is interpreted
as underidentification from singleton evidence, not as a claim that the stack is
unsafe. Pairwise overlap evidence can narrow the interval and show exactly what
additional dependence measurement buys.


\section{Related Work and Limitations}
The mathematical infrastructure is classical. Frechet-Hoeffding bounds,
Boole--Bonferroni inequalities, copulas, and finite coupling arguments long
predate this work \cite{frechet1951,hoeffding1940,boole1854,bonferroni1936,sklar1959,joe1997,nelsen2006}.
Linear programming over partially identified finite distributions is standard
in econometrics and causal inference \cite{manski1990,balke1994,tamer2010,beresteanu2008}.
Reliability engineering also studies common-cause and dependent component
failures \cite{mosleh1998}.

The AI-safety contribution is the translation: composed guardrail evaluation is
cast as a partial-identification problem with explicit failure events,
assumptions, diagnostics, and witnesses. Empirical guardrail failures and
adversarial transfer are well documented in the language-model literature
\cite{wei2023,zou2023,perez2022}, and recent security analyses emphasize that
finite guardrail sets should not be treated as universal robustness mechanisms
\cite{vassilev2026}. HarmBench and JailbreakBench provide public benchmark
contexts for worked examples \cite{mazeika2024harmbench,chao2024jailbreakbench}.

The limitations are central. The binary reduction discards score-level
dependence. The LP output is conditional on supplied evidence. Endpoint
witnesses verify mathematical feasibility, not empirical truth. The explicit
atom LP scales as $2^m$ in the general case. Finally, no interval produced here
is a deployment certificate; risk decisions require external standards,
monitoring, governance, and human accountability.


\section{Conclusion}
Composed guardrail evaluation should distinguish what is observed, what is
assumed, and what follows from those assumptions. The finite atom LP provides a
small but auditable core: it computes sharp bounds for declared binary events,
exposes witness distributions, and makes independence assumptions explicit
rather than implicit. Wide intervals are not a failure of the method; they are
evidence that the available evaluation data underidentify the composed risk.

\bibliographystyle{plain}
\bibliography{references}

\end{document}
